% Source: PDF pages 20--26.
\section{Binary strings and regular expressions}
We have seen how to count all binary strings: $b_n=2^n$ and
\[
B(x)=\sum_{n\ge0}b_nx^n=\sum_{n\ge0}2^nx^n=\frac1{1-2x}.
\]
We now consider binary strings with restrictions.

\subsection{Concatenation}
We write strings as sequences of digits without commas: $(0,1,1,0)=0110$.
\begin{definition}
The \textit{concatenation} of $\sigma=\sigma_1\sigma_2\cdots\sigma_r$ and $\tau=\tau_1\tau_2\cdots\tau_s$ is
\[
\sigma\tau=\sigma_1\sigma_2\cdots\sigma_r\tau_1\tau_2\cdots\tau_s.
\]
\end{definition}
\begin{itemize}
\item The empty string $\varepsilon$ has size $0$ and satisfies $\varepsilon\sigma=\sigma=\sigma\varepsilon$ for every string $\sigma$.
\item We write $\sigma^k=\underbrace{\sigma\sigma\cdots\sigma}_{k\text{ times}}$.
\item We say that $\sigma$ is a \textit{substring} of $\tau$ if $\sigma$ occurs within $\tau$ with no gaps.
\item The \textit{concatenation product} of sets $S$ and $T$ of strings is $ST=\{\sigma\tau:\sigma\in S,\ \tau\in T\}$.
\item We write $S^k=\underbrace{SS\cdots S}_{k\text{ times}}$, $S^0=\{\varepsilon\}$, and $\displaystyle S^*=\bigcup_{k\ge0}S^k$.
\end{itemize}
\begin{remark}
Concatenation can identify different pairs. For example,
\[
\{1,10\}\{1,01\}=\{11,101,101,1001\}=\{11,101,1001\},
\]
whereas the Cartesian product consists of the four distinct pairs
\[
\{(1,1),(10,1),(1,01),(10,01)\}.
\]
\end{remark}
\begin{example}
\begin{itemize}
\item $\{0,1\}^*$ is the set of all binary strings.
\item $\{0,1\}\{0,1\}^*$ is the set of all nonempty binary strings.
\item $(\{0,1\}^2)^*$ is the set of all binary strings of even length.
\end{itemize}
\end{example}
A set of strings can be specified in more than one way. For example,
\[
\{0,1\}^*=\{\varepsilon,0,1\}^*
\]
as sets of strings, but the expressions do not describe the same ways of producing a string: empty strings can be inserted in the second description. We need to consider how strings are produced, not just which strings are obtained.

\subsection{Regular expressions and ambiguity}
\begin{definition}
A \textit{regular expression} is defined recursively by one of the following:
\begin{enumerate}
\item $\varepsilon$, $0$, or $1$;
\item $R\cup S$, where $R$ and $S$ are regular expressions;
\item $RS$, where $R$ and $S$ are regular expressions;
\item $R^*$, where $R$ is a regular expression.
\end{enumerate}
A regular expression $R$ produces a set of strings $\mathcal L(R)$, called a \textit{regular language}.
\end{definition}
We use $0\cup1$ as shorthand for $\{0,1\}$.
\Needspace{11\baselineskip}
\begin{example}
\leavevmode\par
\begin{itemize}
\item $(00)(1\cup01)$ produces $\{001,0001\}$.
\item $(10)^5(111)$ produces $\{1010101010111\}$.
\item $(0\cup1)^*$ produces all binary strings.
\end{itemize}
\end{example}
Regular languages are important in computer science because they are exactly the languages accepted by finite automata. For example, the following automaton accepts all binary strings with no two consecutive $1$s. Double circles indicate accepting states.
\begin{center}
\begin{tikzpicture}[>=Stealth,auto,node distance=3cm,
state/.style={circle,draw,minimum size=8mm},accept/.style={state,double,double distance=1.5pt}]
\node[accept] (a) {$A$};
\node[accept,right=of a] (b) {$B$};
\node[state,right=of b] (c) {$C$};
\draw[->] (-1,.7)--(a);
\path[->] (a) edge[loop below] node {$0$} (a)
 (a) edge[bend left=22] node {$1$} (b)
 (b) edge[bend left=22] node {$0$} (a)
 (b) edge node {$1$} (c)
 (c) edge[loop above] node {$0$} (c)
 (c) edge[loop below] node {$1$} (c);
\end{tikzpicture}
\end{center}

\begin{example}
The expression $(1\cup10)(1\cup01)$ produces
\[
\{1,10\}\{1,01\}=\{11,101,101,1001\}.
\]
The string $101$ is produced in two different ways.
\end{example}
\begin{definition}
A regular expression is \textit{unambiguous} if it produces each string in only one way; otherwise it is \textit{ambiguous}.
\end{definition}
\begin{example}
The expression $(0\cup1)^*$ is unambiguous. The expression $(\varepsilon\cup0\cup1)^*$ is ambiguous, since $1=\varepsilon1=1\varepsilon=\cdots$.
\end{example}
To show that an expression is unambiguous, use the following rules:
\begin{itemize}
\item $\varepsilon$, $0$, and $1$ are unambiguous.
\item Suppose $R$ and $S$ are unambiguous and produce the languages $\mathcal R$ and $\mathcal S$.
Then $R\cup S$ is unambiguous if and only if $\mathcal R\cap\mathcal S=\varnothing$.
\item Under the same assumption, $RS$ is unambiguous if and only if every string in $\mathcal R\mathcal S$ has a unique factorization $rs$ with $r\in\mathcal R$ and $s\in\mathcal S$.
\item $R^*$ is unambiguous if and only if each $R^k$ is unambiguous and the languages $\mathcal R^k$ are pairwise disjoint.
\end{itemize}
\begin{lemma}
The expression $0^*(10^*)^*$ produces all binary strings and is unambiguous.
\end{lemma}
\begin{proof}
For example, consider the following string, broken at its $1$s:
\begin{center}
\begin{tikzpicture}[x=.55cm,y=.65cm]
\foreach \x/\s in {0/0,1/0,2/1,3/0,4/1,5/1,6/0,7/1}{\node at (\x,0) {$\s$};}
\foreach \l/\r in {-.45/1.45,1.55/3.45,3.55/4.45,4.55/6.45,6.55/7.45}{
\draw[notelink,rounded corners=2pt] (\l,-.4) rectangle (\r,.4);
}
\end{tikzpicture}
\end{center}

A string begins with a possibly empty sequence of $0$s, then breaks into pieces beginning with $1$. We can uniquely decompose it using the locations of its $1$s as
\[
\underbrace{(0\cdots0)}_{m_0}
\underbrace{(10\cdots0)}_{1+m_1}\cdots
\underbrace{(10\cdots0)}_{1+m_r},
\]
for some $r\ge0$ and $m_0,m_1,\ldots,m_r\ge0$.
Ranging over all these values gives every binary string, with a unique decomposition.
\end{proof}
\begin{corollary}
The expression $1^*(01^*)^*$ unambiguously produces all binary strings.
\end{corollary}
These expressions can be modified to produce strings with restrictions.
\begin{example}
Give an unambiguous regular expression for the binary strings in which every $0$ is followed by at least one $1$.\par\medskip
Start with $1^*(01^*)^*$. In this expression a $0$ may be followed by any number of $1$s, including none. Requiring at least one gives
\[
1^*(011^*)^*.
\]
This expression is unambiguous and has generating function
\[
\frac1{1-x}\cdot\frac1{1-\frac{x^2}{1-x}}.
\]
\end{example}

Unambiguous expressions translate into generating functions:
\begin{center}
\begin{tabular}{c|c}
Expression & Generating function \\
\hline
$\varepsilon$ & $1$ \\
$0$ or $1$ & $x$ \\
$R\cup S$ & $R(x)+S(x)$ \\
$RS$ & $R(x)S(x)$ \\
$R^*$ & $\displaystyle\frac1{1-R(x)}$
\end{tabular}
\end{center}
\begin{remark}
Unambiguity is the key: union must not count a string twice, concatenation must correspond to the Cartesian product without collisions, and an unambiguous starred expression cannot have the empty string in its repeated language.
\end{remark}

\subsection{Blocks and restrictions}
\begin{definition}
A \textit{block} of a binary string is a nonempty maximal consecutive sequence of the same digit.
\end{definition}
\begin{example}
The string $00011101100$ has blocks
\[
\underbrace{000}_{3}\ \underbrace{111}_{3}\ \underbrace{0}_{1}\
\underbrace{11}_{2}\ \underbrace{00}_{2}.
\]
\end{example}
\begin{theorem}
The expressions
\[
0^*(11^*00^*)^*1^*
\qquad\text{and}\qquad
1^*(00^*11^*)^*0^*
\]
unambiguously produce all binary strings.
\end{theorem}
The idea is to separate the initial possibly empty block, the alternating nonempty blocks, and the final possibly empty block. For example, the first expression starts with $0$s (or nothing), then has pairs of a $1$-block and a $0$-block, and can end in $1$s. The maximal blocks determine the decomposition uniquely.
\begin{example}
Find the generating function for binary strings in which every block has length at most $2$.\par\medskip
Modify $1^*(00^*11^*)^*0^*$ by restricting each block to obtain the unambiguous expression
\[
(\varepsilon\cup1\cup11)((0\cup00)(1\cup11))^*(\varepsilon\cup0\cup00).
\]
Its generating function is
\begin{align*}
B(x)&=(1+x+x^2)\frac1{1-(x+x^2)^2}(1+x+x^2)\\
&=\frac{(1+x+x^2)^2}{(1-x-x^2)(1+x+x^2)}\\
&=\frac{1+x+x^2}{1-x-x^2}.
\end{align*}
\end{example}
\begin{example}
Find the generating function for binary strings in which every block of $0$s has even length.\par\medskip
A nonempty even-length block of $0$s is represented unambiguously by $(00)(00)^*$. Thus the required strings are produced by
\[
1^*\bigl((00)(00)^*11^*\bigr)^*(00)^*.
\]
The generating function is
\begin{align*}
S(x)&=\frac1{1-x}\cdot
\frac1{1-\frac{x^3}{(1-x)(1-x^2)}}\cdot\frac1{1-x^2}\\
&=\frac1{1-x-x^2}.
\end{align*}
\end{example}
\begin{example}
Find the generating function for binary strings that do not contain $011$ as a substring.\par\medskip
Starting with $1^*(00^*11^*)^*0^*$, restrict the $1$-block after any $0$-block to have length at most one. This gives
\[
1^*(00^*1)^*0^*.
\]
The generating function is
\[
\frac1{1-x}\cdot\frac1{1-\frac{x^2}{1-x}}\cdot\frac1{1-x}
=\frac1{1-2x+x^3}.
\]
\end{example}
